\chapter{CONCLUSION}

This project delivered a working statically-typed programming language with the
required primitive types: Integer, Float, Boolean, and String. The language
supports arithmetic expressions, equality and inequality checks, assignments,
if-then-else statements, while-loops, function abstraction with value parameter
passing, return statements, and the built-in \texttt{print()} function.

The final system follows a clean compiler-style architecture. The lexer creates
tokens, the parser builds the AST, the type checker performs static inference
and validation, and the interpreter executes the program. This staged design
made it easier to isolate syntax errors, type errors, and runtime behaviour.

One of the main outcomes of the project is the successful implementation of
static typing without explicit declarations. Variable types are inferred from
assignments, and function signatures are inferred from calls and return
statements. Invalid programs are rejected before runtime. In addition, the
system includes both a command-line runner and a desktop UI so that scripts can
be entered, loaded, executed, and inspected conveniently.

The automated tests added for the semantic and runtime stage pass successfully and confirm correct
behaviour for loop execution, assignment mismatch detection, and function type
inference. These checks provide a baseline for future extensions.

Future work could include recursive functions, richer data structures,
additional operators, and code generation to a lower-level target. However, the
current implementation already satisfies the main functional requirements of the
assignment.
